\documentclass[11pt]{article}
\usepackage{amsmath, amssymb, amsthm}
\usepackage{mathpartir}
\usepackage{listingsutf8}
\usepackage{hyperref}
\usepackage{xcolor}

%%%%%%%%%%%%%%%%%%%%%%%%%%%%%%%%%%%%%%%%%%%%%%%%%%%%%%%%%%%%%%%%%%%%%%%%%%%%%

\newtheorem{definition}{Definition}[section]

% Custom style for code listings
\lstdefinestyle{listing_style}{
    basicstyle=\ttfamily\footnotesize,
    keywordstyle=\color{blue},
    commentstyle=\color{gray},
    stringstyle=\color{gray},
    numbers=left,
    numberstyle=\tiny\color{gray},
    stepnumber=1,
    numbersep=5pt,
    backgroundcolor=\color{white},
    showspaces=false,
    showstringspaces=false,
    showtabs=false,
    tabsize=2,
    captionpos=b,
    breaklines=true,
    breakatwhitespace=false,
    escapeinside={(/}{/)}
}

\lstset{style=listing_style}

\lstdefinelanguage{Ibis}{
    keywords={if, then, else, for, match, def, let, in, site, where, covers, paths, struct, enum},
    sensitive=true,
    morecomment=[l]{--},
    morestring=[b]",
}

%%%%%%%%%%%%%%%%%%%%%%%%%%%%%%%%%%%%%%%%%%%%%%%%%%%%%%%%%%%%%%%%%%%%%%%%%%%%%

\title{Ibis: A purely functional programming language to experiment with categorical memory safety}
\author{Charlotte W.}
\date{\today}

\begin{document}
\maketitle

\tableofcontents

\section{Introduction}
Ibis is an experimental programming language designed to explore the application of category theory to memory safety, and to answer the question: "can we use category theory to reason about memory safety in a systems language?"

Memory regions in Ibis are treated as open sets in a topological space, forming a presheaf over the category of memory regions $\textit{Psh}(\textit{Mem})$. This allows us to reason about memory safety in a categorical framework,
where inclusion morphisms represent safe transformations between memory regions.

\pagebreak

\subsection{Syntax Overview}

\subsection{Site Declarations}
Site declarations are used to define a topological site in Ibis.
Each site declaration consists of a name, a list of covers (open sets), and a list of paths between these functions.

Site declarations inherently map to a cubical type, where the covers are the faces of the cube and the paths are the edges of the cube.

\begin{lstlisting}[language=Ibis, caption={Example of a site declaration in Ibis}]
site MySite where
  covers: [~a, ~b]
  paths:
    from_a_to_b: ~a -> ~b

-- (/$\sigma$/) is the implicit site parameter
def alloc_in {(/$\sigma$/) : MySite} (x: Int) : Ptr Int {(/$\sigma$/)} := ...
\end{lstlisting}

\pagebreak

\section{Category Theory}

\subsection{Yoneda Lemma}
The Yoneda lemma is a foundational result in category theory that states that an object
in a category is defined entirely by its arrows, and how it relates to other objects in the category.

For any locally small category $\mathcal{C}$, an object $X \in \mathcal{C}$, and a presheaf functor $F: \mathcal{C}^{\text{op}} \to \textbf{Set}$, there is a natural bijection:
\[
    \text{Nat}(\text{Hom}_{\mathcal{C}}(-, X), F) \cong F(X)
\]

In other words, natural transformations from the representable hom-functor $\text{Hom}_{\mathcal{C}}(-, X)$ to $F$ are in a one-to-one correspondence with the elements of the set $F(X)$.

This yields the Yoneda Embedding $\textbf{y}: \mathcal{C} \to \textbf{Set}^{\mathcal{C}^{\text{op}}}$, sending $X \mapsto \text{Hom}_{\mathcal{C}}(-, X)$, which fully and faithfully embeds $\mathcal{C}$ into its presheaf category while preserving its internal structure.
\subsection{Opposite Categories and Presheaves}
A presheaf $\mathcal{F}$ over a category $\mathcal{C}$ is a contravariant functor $\mathcal{F} : \mathcal{C}^\text{op} \to \textbf{Set}$, where $\mathcal{C}^\text{op}$ is the opposite category of $\mathcal{C}$ and $\textbf{Set}$ is the category of sets, such that for every morphism $f: X \to Y$ in $\mathcal{C}$, there is a corresponding morphism $f^{\text{op}}: Y \to X$ in $\mathcal{C}^{\text{op}}$.


\subsection{Restriction Maps}
Given a valid morphism $f : \sim\!\text{a} \to \sim\!\textrm{b}$ between covers $\sim\!\text{a}$ and $\sim\!\text{b}$ the restriction map is defined as:

\[
    \text{res}_{U, V} : F(\sim\!\text{a}) \longrightarrow F(\sim\!\text{b})
\]

Morphism reachability is treated as a restriction over a presheaf -- if no valid morphisms exist between two covers, then the restriction map is undefined and it is mathematically impossible to reach one cover from the other.

\subsection{Kan Extensions}
Kan extensions are a fundamental concept in category theory that allow us to extend functors along a given functor.

Given a valid morphism $f : \sim\!\text{a} \to \sim\!\textrm{b}$ between covers $\sim\!\text{a}$ and $\sim\!\text{b}$, the left Kan extension of a functor $F$ along $f$ is defined as:

\[
    \text{ext}_{\sim\!\text{a} \to \sim\!\text{b}} : F(\sim\!\text{a}) \longrightarrow F(\sim\!\text{b}) \quad \text{or} \quad \text{Lan}_f F(\sim\!\text{b})
\]

\begin{definition}[Left Kan Extension]
    Given a functor $F: \mathcal{C} \to \mathcal{E}$ and a mapping $K: \mathcal{C} \to \mathcal{D}$, the Left Kan extension of $F$ along $K$ is a functor $\text{Lan}_K F: \mathcal{D} \to \mathcal{E}$ defined by the colimit for each object $d$ in $\mathcal{D}$, where the colimit is taken over the comma category $(K \downarrow d)$.
    \[
        (\text{Lan}_K F)(d) = \varinjlim_{(b, K(b) \to d)} F(b)
    \]
\end{definition}

In other words, to evaluate the left Kan extension of $F$ along $K$ at an object $d$, we gather all objects $b$ in $\mathcal{C}$ such that there exists
a morphism from $K(b)$ to $d$ in $\mathcal{D}$, and then glue them together using a colimit.

\begin{definition}[Right Kan Extension]
    Given a functor $F: \mathcal{C} \to \mathcal{E}$ and a mapping $K: \mathcal{C} \to \mathcal{D}$, the Right Kan extension of $F$ along $K$ is a functor $\text{Ran}_K F: \mathcal{D} \to \mathcal{E}$ defined by the limit for each object $d$ in $\mathcal{D}$, where the limit is taken over the comma category $(d \downarrow K)$.
    \[
        (\text{Ran}_K F)(d) = \varprojlim_{(d \to K(b), b)} F(b)
    \]
\end{definition}


\subsection{Geometric Morphisms}
A geometric morphism $f: \mathcal{E} \to \mathcal{F}$ between topoi $\mathcal{E}$ and $\mathcal{F}$ consists of a pair of adjoint functors $(f^*, f_*)$ where $f^* \dashv f_*$:

Geometric morphisms are functorial mappings between topoi that preserve the structure of the topos, including finite limits and colimits.

\begin{itemize}
    \item The \textbf{inverse image functor} $f^*: \mathcal{F} \to \mathcal{E}$ is the \emph{left adjoint}. It is left-exact (preserves finite limits) and preserves all colimits, used to pull back contexts and sections.
    \item The \textbf{direct image functor} $f_*: \mathcal{E} \to \mathcal{F}$ is the \emph{right adjoint}. It preserves all limits, used to push forward sections from one site to another.
\end{itemize}

To calculate the direct image $f_*(P)$ of a presheaf $P \in \mathcal{E}$:
\begin{enumerate}
    \item Identify the continuous site functor $u : \mathcal{C}_\mathcal{F} \to \mathcal{C}_\mathcal{E}$ associated with the geometric morphism $f$.
    \item For each object $V \in \mathcal{C}_\mathcal{F}$, evaluate $P$ at its preimage object $u(V) \in \mathcal{C}_\mathcal{E}$, yielding $f_*(P)(V) = P(u(V))$.
    \item For each target morphism $g : V' \to V$ in $\mathcal{C}_\mathcal{F}$, map it to the corresponding morphism $u(g) : u(V') \to u(V)$ in $\mathcal{C}_\mathcal{E}$.
    \item Apply the restriction map $\text{res}_{u(V), u(V')}$ to $P(u(V))$ to obtain $P(u(V'))$, ensuring that the presheaf structure is preserved.
\end{enumerate}

To calculate the inverse image $f^*(Q)$ of a sheaf/presheaf $Q \in \mathcal{F}$:
\begin{enumerate}
    \item \textbf{Construct the Comma Category (Index Graph):}
          For a source object $U \in \mathcal{C}_\mathcal{E}$, form the comma category $(U \downarrow u)$ consisting of pairs $(V, \phi)$ where $V \in \mathcal{C}_\mathcal{F}$ and $\phi : U \to u(V)$ is a morphism in $\mathcal{C}_\mathcal{E}$.

    \item \textbf{Evaluate the Filtered Colimit (Raw Pullback $f_p^* Q$):}
          Compute the raw presheaf pullback at $U$ by taking the colimit of section values over the comma category:
          \[
              (f_p^* Q)(U) = \varinjlim_{(V, \phi) \in (U \downarrow u)} Q(V)
          \]
          This collects all candidate sections $Q(V)$ across target contexts that cover $U$, identifying sections that agree under restrictions.

    \item \textbf{First Plus Construction Pass ($(f_p^* Q)^+$):}
          For each object $U$ and covering sieve $R \in J(U)$, gather matching families of sections in $f_p^* Q$ that agree on overlaps to construct local gluing candidates:
          \[
              (f_p^* Q)^+(U) = \varinjlim_{R \in J(U)} \text{Matching}(R, f_p^* Q)
          \]

    \item \textbf{Second Plus Construction Pass ($f^* Q = ((f_p^* Q)^+)^+$):}
          Re-apply the plus construction pass to enforce unique local existence and unicity of glued sections, yielding a true sheaf $f^* Q \in \mathcal{E}$.
\end{enumerate}

\end{document}